\documentclass[answers,addpoints,11pt]{exam}
\usepackage[utf8]{inputenc}
\usepackage[english]{babel}
\usepackage{listings}
\usepackage{amsmath}
\usepackage{algpseudocode}
\usepackage{graphicx}
\usepackage{amssymb}
\usepackage{amsthm}
\usepackage{hyperref}
\usepackage{logicproof}
\usepackage{multicol}
\usepackage{environ}
\usepackage{xcolor}
\usepackage{qtree}

\newcommand{\limn}[1]{ \lim_{n \rightarrow \infty} \left( #1 \right) }
\bracketedpoints
\qformat{Question \thequestion{} : \totalpoints{} \points \hfill}

%These setting will make the code areas look Pretty
\lstset{
	escapechar=~,
	numbers=left,
	%numberstyle=\tiny,
	stepnumber=1,
	firstnumber=1,
	%numbersep=5pt,
	language=C,
	stringstyle=\ttfamily,
	%basicstyle=\footnotesize,
	showstringspaces=false,
	frame=single
}

\title{\textbf{CS260 $|$ Data Structures\\Written Assignment 1}}
\author{\textbf{Srinarayan Srikanth ss5837}}
\date{}

\pagestyle{foot}
\runningfootrule
\footer{}{Page \thepage\ of \numpages}{}

\begin{document}
\maketitle

\section{Questions}
\begin{questions}
%%%%%%%%%%%%%%%%%%%%%%%%%%
%%%% Question 1 %%%%%%%%%%
%%%%%%%%%%%%%%%%%%%%%%%%%%
\question

For each of the following code snippets:

\begin{itemize}
\item Give a tight bound on the runtime of the snippet, and
\item briefly explain your bound.
\end{itemize}

If the runtime is asymptotically different in different cases, give a separate bound for each case and briefly explain what the cases are, e.g. ``insertion sort runs in $\Theta(n)$ time in the best case, when the list is already sorted, because...".


\begin{parts}
\part[7] \texttt{foo}

\begin{verbatim}
/* L is array of n integers */
function foo(L, n)
{
  a = 0
  while (a < n) {
    b = a
    x = 0
    while (b < n) {
      x = x + L[b]
      b++
    }
    L[a] = x
    a++
  }
  return L
}

\end{verbatim}

\begin{solution}
\textbf{Bound:} $\Theta(n^2)$.

The outer \texttt{while} loop for $a$ runs from $0$ to $n-1$. For each value of $a$, the inner \texttt{while} loop for $b$ starts at $a$ and runs up to $n-1$, so it executes $n-a$ times. The total number of inner-loop iterations is
\[
\sum_{a=0}^{n-1} (n-a) = \frac{n(n+1)}{2},
\]
which is $\Theta(n^2)$.
\end{solution}

\part[7] \texttt{baz}

\begin{verbatim}
/* L is a list of n integers */
function baz(L, n) {
  for (i = 0; i < n; i++) {
    if (L[i] == 7) {
      return true
    }
  }
  return false
}
\end{verbatim}

\begin{solution}
\textbf{Best case:} $\Theta(1)$, which occurs when $L[0] = 7$, so the function returns \texttt{true} immediately on the first iteration.

\textbf{Worst case:} $\Theta(n)$, which occurs when $7$ does not appear in $L$ (or appears only at $L[n-1]$), so the loop must run through all $n$ elements before returning.
\end{solution}


\part[7] \texttt{bar}

\begin{verbatim}
/* a, b, c, d are all integers */
function bar(a, b, c, d) {
  if (a < b) {
    if (c < d) {
      if (a < c) {
        return a
      } else {
        return c
      }
    } else {
      if (a < d) {
        return a
      } else {
        return d
      }
    }
  } else {
    if (c < d) {
      if (b < c) {
        return b
      } else {
        return c
      }
    } else {
      if (b < d) {
        return b
      } else {
        return d
      }
    }
  }
}
\end{verbatim}

\begin{solution}
\textbf{Bound:} $\Theta(1)$.

The function is a fixed sequence of nested \texttt{if-else} comparisons with no loops or recursion. Every path through the function performs exactly three comparisons before returning, which is constant work independent of any input size.
\end{solution}


\part[7] \texttt{qux}

\begin{verbatim}
/* n is a positive integer */
function qux(n) {
  for (i = 0; i < n; i++) {
    for (j = 0; j < n; j++) {
      for (k = 0; k < n; k++) {
        for (l = 0; l < 8; l++) {
          print(i * (j + k) - l)
        }
      }
    }
  }
}
\end{verbatim}

\begin{solution}
\textbf{Bound:} $\Theta(n^3)$.

The three outer loops ($i$, $j$, $k$) each run exactly $n$ times, and the innermost loop ($l$) always runs exactly $8$ times (a constant), independent of $n$. The total number of \texttt{print} calls is $8n^3$, which is $\Theta(n^3)$.
\end{solution}


\part[7] \texttt{xyzzy}

\begin{verbatim}
/* n is a positive integer */
function xyzzy(n) {
  x = 2
  for (i = 0; i < n; i++) {
    x = x * 7
  }
  for (j = 0; j < 2*n; j++) {
    x = x / 3
  }
  return x
}

\end{verbatim}

\begin{solution}
\textbf{Bound:} $\Theta(n)$.

The two \texttt{for} loops are sequential, not nested. The first runs $n$ times and the second runs $2n$ times, each doing constant work per iteration. The total is $n + 2n = 3n$, which is $\Theta(n)$.
\end{solution}

\end{parts}

%%%%%%%%%%%%%%%%%%%%%%%%%%
%%%% Question 2 %%%%%%%%%%
%%%%%%%%%%%%%%%%%%%%%%%%%%
\question[20]
    Order the following functions from least to greatest growth rate, such that for any two functions $f$ and $g$, if $f$ comes before $g$, then $f(n) = O(g(n))$:

\begin{itemize}
\item $f_1(n) = n^n$
\item $f_2(n) = n^{0.00001}$
\item $f_3(n) = n^2$
\item $f_4(n) = \lg(n)$
\item $f_5(n) = 2^n$
\end{itemize}

\textbf{You must prove each ranking}, but you may assume the transitive property is true, which says that if $f(n) = O(g(n))$ and $g(n) = O(h(n))$, then $f(n) = O(h(n))$.

This means that if the functions were $f, g, h$ and you gave the order $f, g, h$, then you'd only need to prove that $f(n) = O(g(n))$ and $g(n)= O(h(n))$, as you may assume that this implies $f(n) = O(h(n))$.

\textbf{You must fully simplify your limit ratios to receive credit.} If you haven't simplified your limit ratio to 0, $\infty$, or a positive constant, you're not finished simplifying.


\begin{solution}
\textbf{Order (least to greatest):} $f_4(n),\ f_2(n),\ f_3(n),\ f_5(n),\ f_1(n)$

\medskip
\textbf{1. $f_4(n) = O(f_2(n))$:}
\[
\limn{ \frac{\lg(n)}{n^{0.00001}} } = \limn{ \frac{\ln(n)/\ln 2}{n^{0.00001}} }
\]
Applying L'H\^opital's rule (both numerator and denominator $\to \infty$):
\[
= \limn{ \frac{1/(n\ln 2)}{0.00001\, n^{-0.99999}} } = \limn{ \frac{1}{0.00001 \ln 2 \cdot n^{0.00001}} } = 0.
\]
So $f_4(n) = O(f_2(n))$.

\medskip
\textbf{2. $f_2(n) = O(f_3(n))$:}
\[
\limn{ \frac{n^{0.00001}}{n^2} } = \limn{ n^{-1.99999} } = 0.
\]
So $f_2(n) = O(f_3(n))$.

\medskip
\textbf{3. $f_3(n) = O(f_5(n))$:}
\[
\limn{ \frac{n^2}{2^n} }
\]
Applying L'H\^opital's rule twice:
\[
= \limn{ \frac{2n}{2^n \ln 2} } = \limn{ \frac{2}{2^n (\ln 2)^2} } = 0.
\]
So $f_3(n) = O(f_5(n))$.

\medskip
\textbf{4. $f_5(n) = O(f_1(n))$:}
\[
\limn{ \frac{2^n}{n^n} } = \limn{ \left(\frac{2}{n}\right)^n } = 0,
\]
since $2/n \to 0$ and is raised to an ever-growing power. So $f_5(n) = O(f_1(n))$.

\medskip
By transitivity, $f_4(n) = O(f_2(n)) = O(f_3(n)) = O(f_5(n)) = O(f_1(n))$, giving the order above.
\end{solution}

%%%%%%%%%%%%%%%%%%%%%%%%%%
%%%% Question 3 %%%%%%%%%%
%%%%%%%%%%%%%%%%%%%%%%%%%%
\question
Consider the following function \texttt{compute}:

\begin{lstlisting}
/* n is a positive integer */
function compute(n)
{
  if (n == 0)
  {
    return 0
  }
  x = 0
  for (i = 1; i <= n; i++)
  {
    x = x + i
  }
  return compute(n - 1) + x
}
\end{lstlisting}

\begin{parts}

%

\part[5] Give a recurrence that represents the runtime of \texttt{compute}.

\begin{solution}
The \texttt{for} loop does $\Theta(n)$ work, followed by one recursive call on $n-1$:
\[
T(n) = T(n-1) + \Theta(n), \qquad T(0) = \Theta(1).
\]
\end{solution}

%

\part[10] Solve your recurrence using the substitution method, then give a tight bound on the runtime of \texttt{compute}.

\begin{solution}
Write $T(n) = T(n-1) + n$. Expanding:
\[
T(n) = T(n-1) + n = T(n-2) + (n-1) + n = T(n-3) + (n-2) + (n-1) + n = \cdots
\]
After $k$ substitutions:
\[
T(n) = T(n-k) + \sum_{i=0}^{k-1} (n-i).
\]
The base case occurs when $n - k = 0$, i.e. $k = n$:
\[
T(n) = T(0) + \sum_{i=0}^{n-1} (n-i) = T(0) + \sum_{j=1}^{n} j = \Theta(1) + \frac{n(n+1)}{2}.
\]
\textbf{Tight bound:} $T(n) = \Theta(n^2)$.
\end{solution}

\end{parts}

%%%%%%%%%%%%%%%%%%%%%%%%%%
%%%% Question 4 %%%%%%%%%%
%%%%%%%%%%%%%%%%%%%%%%%%%%
\question
Professor Holmes has come up with a new sorting algorithm. He calls it \textbf{Trimerge Sort} and $``$claims$"$ that it is asymptotically faster than \textbf{Merge Sort}, despite the fact both the algorithms operate using similar logic. But, unlike Merge Sort, Trimerge Sort splits the input array into \textit{(roughly)} 3 equal parts at each step of the recursion as long as the array is splittable \textit{(i.e., has at least 3 elements)}. Trimerge Sort's merge subroutine, similar in principle to the one used by Merge Sort, takes 3 sorted subarrays and merges them to produce a single sorted array. Given all of this, answer the following questions.

\begin{parts}
 \part[5] In Merge Sort, the \textbf{merge subroutine} makes $n$ comparisons to merge $2$ arrays of size $n/2$, which takes $\Theta(n)$ time. How many comparisons (in exact numbers) will the merge subroutine of Trimerge Sort make to merge $3$ arrays of size $n/3$? What would be the $\Theta(\cdots)$ bound on the running time for merge subroutine?

 \begin{solution}
 To place each of the $n$ elements into the merged output, the subroutine must find the minimum of the current front elements of the 3 subarrays. Finding the minimum of 3 elements takes 2 comparisons (compare the first two, then compare that winner against the third). Doing this for each of the $n$ elements placed gives
\[
2n \text{ comparisons},
\]
matching the same style of count given for the 2-array case ($n$ comparisons for 2 arrays $\Rightarrow$ $2n$ for 3 arrays, since one extra comparison is needed per element to pick among 3 candidates instead of 2).

\textbf{Bound:} $\Theta(n)$.
 \end{solution}

\part[5] What is the recurrence relation for the Trimerge Sort algorithm? Also explain it in plain English what each term in the recurrence relation stands for.
 \begin{solution}
\[
T(n) = 3T(n/3) + \Theta(n)
\]
\begin{itemize}
\item $T(n)$: total time to sort an array of size $n$ with Trimerge Sort.
\item $3T(n/3)$: the cost of recursively sorting the 3 subarrays, each of (roughly) size $n/3$.
\item $\Theta(n)$: the cost of the merge subroutine that combines the 3 sorted subarrays back into one sorted array of size $n$.
\end{itemize}
 \end{solution}

 \part[8] Prove the $\Theta(\cdots)$ bound on the running time of the Trimerge Sort algorithm using the \textbf{Recursion Tree Method} covered in class.
 \begin{solution}
Let the merge cost at a node handling $m$ elements be $cm$.

\begin{itemize}
\item \textbf{Level 0 (root):} 1 node of size $n$, cost $cn$.
\item \textbf{Level 1:} 3 nodes of size $n/3$, total cost $3 \cdot c(n/3) = cn$.
\item \textbf{Level $i$:} $3^i$ nodes, each of size $n/3^i$, total cost $3^i \cdot c(n/3^i) = cn$.
\item \textbf{Height:} the array shrinks by a factor of 3 each level until reaching size 1, so the tree has $\log_3 n$ levels.
\end{itemize}

Every level costs $cn$, and there are $\log_3 n$ levels, so the total work is
\[
\sum_{i=0}^{\log_3 n} cn = cn\log_3 n.
\]
\textbf{Tight bound:} $\Theta(n \log n)$.
 \end{solution}

  \part[8] Prove the $\Theta(\cdots)$ bound on the running time of the Trimerge Sort algorithm using \textbf{Equation Expansion} method over the recurrence relation you came up with in part (b).
 \begin{solution}
Starting from $T(n) = 3T(n/3) + cn$, expand repeatedly:
\[
T(n) = 3\left(3T(n/9) + c\frac{n}{3}\right) + cn = 9T(n/9) + 2cn
\]
\[
T(n) = 9\left(3T(n/27) + c\frac{n}{9}\right) + 2cn = 27T(n/27) + 3cn
\]
After $k$ expansions:
\[
T(n) = 3^k T(n/3^k) + kcn.
\]
The recursion bottoms out when $n/3^k = 1$, i.e. $k = \log_3 n$. Substituting:
\[
T(n) = 3^{\log_3 n} T(1) + (\log_3 n)\, cn = n\cdot\Theta(1) + cn\log_3 n.
\]
\textbf{Tight bound:} $\Theta(n \log n)$.
 \end{solution}

 \part[4] Is Trimerge Sort asymptotically faster than Merge Sort? Why/why not? Reason through your answer using \textbf{a mathematical proo}f.

 \textit{(\textbf{Hint:} Use the limit method demonstrated in class as well as the recitation on Asymptotic Analysis $\&$ Recurrence Relations)}
 \begin{solution}
Merge Sort's recurrence, $T_m(n) = 2T_m(n/2) + \Theta(n)$, solves to $\Theta(n\log_2 n)$. Trimerge Sort solves to $\Theta(n\log_3 n)$ (parts (c), (d)). Compare via the limit ratio:
\[
\limn{ \frac{n\log_3 n}{n\log_2 n} } = \limn{ \frac{\ln n / \ln 3}{\ln n / \ln 2} } = \frac{\ln 2}{\ln 3} \approx 0.63,
\]
a positive constant. Since the limit is a positive, finite constant, $n\log_3 n = \Theta(n\log_2 n)$, so the two functions have the \emph{same} asymptotic growth rate.

\textbf{Conclusion:} Trimerge Sort is \emph{not} asymptotically faster than Merge Sort; both run in $\Theta(n\log n)$ time.
 \end{solution}

\end{parts}

\end{questions}
\end{document}
